\subsection{The set of threshold solutions}

Under the preceding smooth semipositive dual hypothesis, the divisor relation also has a converse. Once one solution is fixed, positivity is the only remaining condition on its divisorial coefficients.

\begin{proposition}\label{prop:threshold-solution-polytope}
Assume the hypotheses of \cref{thm:semipositive-dual-exact-equation}. Fix a threshold solution $u_0$, and write $Z_1,\ldots,Z_m$ for the prime divisorial components of $Z$. Put $\lambda_j=\nu(u_0,Z_j)$ and
\[
 \mathcal P=\left\{r\in\mathbb R_{\geq0}^m:
       \sum_{j=1}^m r_j\{Z_j\}
          =\sum_{j=1}^m\lambda_j\{Z_j\}
           \text{ in }H^{1,1}_{\mathrm{BC}}(X,\mathbb R)\right\}.
\]
Then $\mathcal P$ is a nonempty compact convex polytope. The map
\[
 [u]\longmapsto\bigl(\nu(u,Z_1),\ldots,\nu(u,Z_m)\bigr)
\]
is a bijection from all threshold subsolutions modulo constants to $\mathcal P$. All these subsolutions solve the exact equation. After normalization by zero integral against a fixed Gauduchon volume, the inverse map is affine and continuous into $L^1$; it is a homeomorphism onto the solution set with its $L^1$ topology.

The residual current
\[
 R=\omega+\ddc u-\sum_{j=1}^m\nu(u,Z_j)[Z_j]
\]
is positive and independent of the solution. In particular uniqueness modulo constants is equivalent to $\mathcal P$ being a singleton. Distinct points of $\mathcal P$ give distinct multiplier-ideal classes and distinct ordinary singularity types.
\end{proposition}
\begin{proof}
The preceding theorem shows that every threshold subsolution is an exact solution and that
\[
 \ddc(u-u_0)=\sum_j\bigl(\nu(u,Z_j)-\lambda_j\bigr)[Z_j].
\]
Its class relation puts the coefficient vector in $\mathcal P$. It also proves injectivity modulo constants and independence of $R$. Positivity of $R$ follows from \cref{lem:nonclosed-divisorial-removal} applied to $u_0$.

Fix a Gauduchon metric $\gamma$ and use the linear decomposition of \cref{lem:constant-trace-bc-representative}. Write
\[
 [Z_j]=B_j+\ddc\ell_j,\qquad
 B_j=B_\gamma([Z_j]),\qquad \int_X\ell_j\gamma^n=0.
\]
The $B_j$ are the chosen smooth representatives of the divisor classes. Since this choice is linear on Bott--Chern cohomology, $r\in\mathcal P$ implies $\sum_j(r_j-\lambda_j)B_j=0$. Normalize $u_0$ to have zero integral and set, initially as an $L^1$ distribution,
\[
 u_r=u_0+\sum_j(r_j-\lambda_j)\ell_j.
\]
Then
\[
 \omega+\ddc u_r=R+\sum_j r_j[Z_j]\geq0.
\]
Thus $u_r$ has a genuine $\omega$-plurisubharmonic representative. The residual current has zero generic coefficient along each $Z_j$, so $\nu(u_r,Z_j)=r_j$.

On $X\setminus Z$ the difference $u_r-u_0$ has zero complex Hessian and is smooth by elliptic regularity. Hence the smooth Monge--Amp\`ere measures of $u_r$ and $u_0$ agree there. Their extended measures give $Z$ no mass. It follows that $\mathrm{MA}^{\mathrm{ext}}_\omega(u_r)=bf\,\mathrm{d}V$ on $X$, proving surjectivity. The displayed formula for $u_r$ is affine and continuous in $L^1$ and preserves the zero-integral normalization.

The set $\mathcal P$ contains $\lambda$ and is the intersection of an affine subspace with finitely many closed half-spaces. To prove boundedness, put $d_j=\int_{Z_j}\gamma^{n-1}>0$. Pairing its class relation with $\gamma^{n-1}$ gives
\[
 \sum_j d_jr_j=\sum_j d_j\lambda_j.
\]
Since $r_j\geq0$, every coordinate is bounded. Thus $\mathcal P$ is compact, and its continuous bijection to the normalized solution set in the Hausdorff space $L^1$ is a homeomorphism. If $m=0$, the coefficient space is a singleton and the same conclusion follows directly. Finally, equality of multiplier ideals forces equality of these generic Lelong numbers by \cref{thm:local-i-inputs}; ordinary equivalence implies multiplier-ideal equivalence.
\end{proof}

This description concerns solutions for the fixed smooth semipositive maximizing form and its analytic degeneracy set. It gives no criterion for existence of such a form and does not classify threshold solutions when only singular current attainment is known.
